\chapter{Synchronous Models}

\section{Core assumption}
\begin{defbox}
A synchronous system assumes known bounds on message delays, process speeds, and clock drift. These bounds make timing part of the model rather than a source of uncertainty.
\end{defbox}

\section{Typical properties}
\begin{itemize}
  \item Messages arrive within a known maximum time.
  \item Local computation steps have a known upper bound.
  \item Protocols can use timeouts as reliable evidence of failure or delay.
\end{itemize}

\section{Why it matters}
\begin{keybox}
Many classical distributed algorithms are simpler in a synchronous model because they can distinguish between a slow process and a failed one using timing bounds.
\end{keybox}

\section{Exam prompts}
\begin{itemize}
  \item State the assumptions that define synchrony.
  \item Explain how synchrony supports timeout-based reasoning.
  \item Compare synchrony with asynchronous communication.
\end{itemize}
